\section{Half-line Sturm--Liouville theory as the source of Jost data}
\label{sec:sturm-liouville-jost}
% BEGIN CURATED SUBJECT INDEX
\index{operator!adjoint and self-adjoint}
\index{spectral measure}
\index{Sturm--Liouville theory}
\index{Jost function}
\index{Wronskian}
\index{partial-wave expansion}
\index{boundary condition!incoming and outgoing}
\index{resonance!residue}
\index{exact WKB}
\index{radial equation}
% END CURATED SUBJECT INDEX
\index{Sturm--Liouville operator}\index{Jost function}

\calculationroute
  {The radial Schr\"odinger equation, partial waves, and Jost asymptotics of
  Section~\ref{sec:textbook-scattering}.}
  {Construct the self-adjoint half-line operator from its differential
  expression, derive its Green kernel and spectral measure, and then continue
  the Wronskian that becomes the resonance condition.}
  {Given a radial potential, classify both endpoints, calculate the regular
  and outgoing solutions, form their Wronskian, and differentiate it to
  obtain a pole residue.}

A radial equation becomes a scattering problem only after two logically
different choices have been made.  At the origin one chooses an operator
domain.  At infinity one chooses a comparison solution and, after analytic
continuation, a sheet.  Sturm--Liouville theory controls the first choice and
the physical spectral measure; Jost theory adds the analytic data required
for the second.  Exact WKB joins them by transporting the origin solution to
the asymptotic sector and calculating their Wronskian.

The purpose of this section is not to place an abstract theorem beside the
radial equation.  We derive each object in the order in which it is used in a
calculation.  References for the underlying spectral theory include
Refs.~\cite{Titchmarsh:1962,Zettl:2005,Rakityansky:2022Jost,
Kukulin:1989}.

\subsection{The differential expression and the Hilbert space}
% BEGIN CURATED SUBJECT INDEX
\index{operator!adjoint and self-adjoint}
\index{Hilbert space}
\index{measure theory}
\index{radial equation}
% END CURATED SUBJECT INDEX
\index{maximal operator}\index{minimal operator}\index{boundary form}
\label{subsec:sl-expression-space}

On an interval $(a,b)$, possibly with singular endpoints, consider
\begin{equation}
  \tau u
  =\frac{1}{w(r)}
  \left[-\frac{\rmd}{\rmd r}
  \left(p(r)\frac{\rmd u}{\rmd r}\right)+q(r)u(r)\right],
  \label{eq:sturm-liouville-expression}
\end{equation}
where $p>0$, $w>0$ almost everywhere and
$1/p,q,w$ are locally integrable.  The natural Hilbert space is
\begin{equation}
  \Hil_w=L^2((a,b),w(r)\rmd r),
  \qquad
  \dualpair{u}{v}_w
  =\int_a^b u(r)^*v(r)w(r)\dd r.
  \label{eq:sl-weighted-space}
\end{equation}
The maximal domain consists of functions for which $u$ and $pu'$ are locally
absolutely continuous and $\tau u\in\Hil_w$.  Absolute continuity, rather
than twice continuous differentiability, is enough to make integration by
parts meaningful and includes many potentials used in scattering.

For $u,v$ in the maximal domain define the boundary bracket
\begin{equation}
  [u,v](r)
  =p(r)\left(u(r)^*v'(r)-u'(r)^*v(r)\right).
  \label{eq:sl-boundary-bracket}
\end{equation}
Direct integration gives Lagrange's identity
\begin{equation}
  \dualpair{u}{\tau v}_w
  -\dualpair{\tau u}{v}_w
  =[u,v](b)-[u,v](a),
  \label{eq:lagrange-identity}
\end{equation}
where endpoint values are limits from inside the interval.  Equation
\eqref{eq:lagrange-identity} is the radial version of probability-current
conservation.  It also shows why the differential formula alone does not
define a self-adjoint Hamiltonian: a domain must make the boundary form
vanish for every pair of domain vectors.

For the reduced radial Schr\"odinger equation one has $p=w=1$ after
multiplying by $2\mu/\hbar^2$,
\begin{equation}
  \tau_l u
  =-u''+\left[\frac{l(l+1)}{r^2}+U(r)\right]u,
  \qquad
  U(r)=\frac{2\mu}{\hbar^2}V(r),
  \label{eq:sl-radial-expression}
\end{equation}
and the spectral parameter is $k^2=2\mu E/\hbar^2$.  The reduced function
$u$ belongs to $L^2(0,\infty;\rmd r)$ for a bound state; the factor $r^2$ in
the three-dimensional volume element has already been removed.

\subsection{Regular endpoints and the boundary condition at the origin}
\label{subsec:sl-regular-endpoint}
% BEGIN CURATED SUBJECT INDEX
\index{operator!adjoint and self-adjoint}
\index{self-adjoint extension}
% END CURATED SUBJECT INDEX

An endpoint $a$ is regular if it is finite and $1/p,q,w$ are integrable near
it.  The boundary values $u(a)$ and $(pu')(a)$ then exist.  A separated
self-adjoint condition is
\begin{equation}
  u(a)\cos\alpha+(pu')(a)\sin\alpha=0,
  \qquad \alpha\in[0,\pi).
  \label{eq:sl-separated-boundary}
\end{equation}
Substituting two functions obeying the same condition into
Eq.~\eqref{eq:sl-boundary-bracket} makes the bracket vanish.  Dirichlet and
Neumann conditions are the cases $\alpha=0$ and $\alpha=\pi/2$.

For the three-dimensional $s$ wave with a nonsingular potential, the reduced
operator $-\rmd^2/\rmd r^2$ on $(0,\infty)$ needs a boundary condition at
$r=0$.  Regularity of the full wave function $u(r)/r$ selects
\begin{equation}
  u(0)=0.
  \label{eq:radial-dirichlet-origin}
\end{equation}
Other self-adjoint extensions describe additional contact physics and should
not be introduced accidentally by a numerical boundary stencil.  The phrase
``regular solution'' means the solution satisfying the declared origin
domain, conventionally normalized by
\begin{equation}
  \varphi(k,0)=0,
  \qquad
  \varphi'(k,0)=1
  \label{eq:regular-solution-initial}
\end{equation}
for an $s$ wave.

\subsection{Singular endpoints and Weyl's alternative}
% BEGIN CURATED SUBJECT INDEX
\index{operator!adjoint and self-adjoint}
\index{unbounded operator}
\index{Sturm--Liouville theory!limit-point/limit-circle}
\index{boundary condition!incoming and outgoing}
\index{Langer correction}
\index{Frobenius solution}
% END CURATED SUBJECT INDEX
\index{limit-point endpoint}\index{limit-circle endpoint}
\index{Weyl alternative}
\label{subsec:weyl-limit-point-circle}

At a singular endpoint, individual boundary values may not exist.  Weyl's
classification asks instead how many solutions of
\begin{equation}
  \tau u=z u,
  \qquad z\in\complex\setminus\real,
  \label{eq:sl-nonreal-equation}
\end{equation}
are square integrable near that endpoint:
\begin{itemize}
  \item the endpoint is \textit{limit circle} if every solution is locally in
  $L^2_w$ there;
  \item it is \textit{limit point} if, up to scale, at most one solution is
  locally in $L^2_w$ there.
\end{itemize}
The classification is independent of the chosen nonreal $z$.  A limit-circle
endpoint requires one boundary condition.  At a limit-point endpoint the
boundary bracket vanishes automatically on the operator domain, and no
boundary condition is added.  Imposing one anyway can overconstrain the
problem.

The inverse-square model makes the criterion calculable.  Near $r=0$ let
\begin{equation}
  \tau_\nu=-\frac{\rmd^2}{\rmd r^2}
  +\frac{\nu^2-\tfrac14}{r^2}.
  \label{eq:inverse-square-endpoint-model}
\end{equation}
The two Frobenius behaviors are
\begin{equation}
  u_+(r)\sim r^{1/2+\nu},
  \qquad
  u_-(r)\sim r^{1/2-\nu}
  \label{eq:inverse-square-frobenius}
\end{equation}
for $\nu\ne0$, with a logarithmic replacement at $\nu=0$.  Since
$\int_0^\epsilon r^{1-2\nu}\rmd r$ converges exactly when $\nu<1$, the
origin is limit circle for $0\leq\nu<1$ and limit point for $\nu\geq1$.
For $l(l+1)/r^2$, $\nu=l+1/2$.  Thus the bare $s$-wave endpoint needs the
regular boundary condition, whereas $l\geq1$ is limit point.  The Langer map
of Stage V replaces $l(l+1)$ by $(l+1/2)^2$ in the WKB momentum; it does not
erase this operator-domain analysis.

For a sufficiently short-range potential, $r=\infty$ is limit point.  A
self-adjoint half-line Hamiltonian therefore uses the origin condition but no
independent $L^2$ condition at infinity.  Bound-state square integrability
emerges from the spectral equation.  Outgoing resonance behavior is an
analytic continuation of a comparison solution, not a second self-adjoint
boundary condition on the original operator.

\subsection{Fundamental solutions, the Weyl solution, and the Green kernel}
\label{subsec:weyl-m-green}
% BEGIN CURATED SUBJECT INDEX
\index{spectral measure}
\index{resolvent}
\index{Sturm--Liouville theory!limit-point/limit-circle}
\index{Weyl--Titchmarsh function}
\index{Wronskian}
% END CURATED SUBJECT INDEX
\index{Weyl--Titchmarsh function}\index{Green function}

Choose real entire fundamental solutions $\theta(z,r)$ and $\phi(z,r)$ at a
regular left endpoint so that
\begin{equation}
  W[\theta,\phi]
  :=p(\theta\phi'-\theta'\phi)=1.
  \label{eq:sl-fundamental-wronskian}
\end{equation}
The solution $\phi$ is chosen to obey the left boundary condition.  If the
right endpoint is limit point, then for $z\notin\real$ there is a unique
coefficient $m(z)$ such that
\begin{equation}
  \psi(z,r)=\theta(z,r)+m(z)\phi(z,r)
  \label{eq:weyl-solution}
\end{equation}
is square integrable near $b$.  The scalar $m(z)$ is the Weyl--Titchmarsh
function.  With the conventional choice of fundamental solutions it is a
Herglotz function,
\begin{equation}
  \im z>0
  \quad\Longrightarrow\quad
  \im m(z)>0,
  \label{eq:m-herglotz}
\end{equation}
and has an integral representation whose measure is the spectral measure of
the half-line operator.

In the convention used throughout the book, let
$\Resolvent_H(z)=(z-H)^{-1}$.  If $u_a(z,r)$ obeys the left
boundary condition and $u_b(z,r)$ is the Weyl solution at the right, then
\begin{equation}
  G(z;r,r')
  =\frac{u_a(z,r_<)u_b(z,r_>)}
  {W[u_a,u_b](z)}
  \label{eq:sl-green-kernel}
\end{equation}
solves $(z-H)G=\delta(r-r')/w(r')$.  To check the sign, integrate across
$r=r'$: the jump in $p\partial_rG$ is $+1$, exactly as required by the
leading term $+(pu')'$ in $z-H$.

Equation~\eqref{eq:sl-green-kernel} is more than a formula for inversion.  It
shows immediately that a zero of the Wronskian is a pole of the Green
function and that the numerator factorizes at that zero.  The remaining task
is to identify $u_b$ on the physical continuum and to continue it.

\subsection{The spectral measure and generalized eigenfunction expansion}
\label{subsec:sl-eigenfunction-expansion}
% BEGIN CURATED SUBJECT INDEX
\index{operator!adjoint and self-adjoint}
\index{spectral theorem}
\index{spectral measure}
\index{direct integral}
\index{Gamow--Siegert state}
% END CURATED SUBJECT INDEX

The boundary values of $m$ recover the spectral measure $\rho$.  At an
absolutely continuous energy one has, in the distributional sense,
\begin{equation}
  \frac{\rmd\vsub{\rho}{ac}}{\rmd E}
  =\frac{1}{\pi}\im m(E+\rmi0),
  \label{eq:m-spectral-density}
\end{equation}
while poles of $m$ on the real axis supply point masses.  If $\phi(E,r)$ is
the left-regular solution in the same normalization, the spectral transform
is
\begin{align}
  \widehat f(E)
  &=\int_a^b\phi(E,r)f(r)w(r)\dd r,
  \label{eq:sl-spectral-transform}\\
  f(r)
  &=\int_\real\phi(E,r)\widehat f(E)\dd\rho(E),
  \label{eq:sl-spectral-inverse}\\
  \|f\|_w^2
  &=\int_\real|\widehat f(E)|^2\dd\rho(E).
  \label{eq:sl-parseval}
\end{align}
The continuum normalization is therefore not guessed from a delta function.
It is determined by the pair consisting of the regular solution and its
spectral measure.  Rescaling $\phi(E,r)$ rescales $\rmd\rho(E)$ inversely and
leaves Eqs.~\eqref{eq:sl-spectral-inverse}--\eqref{eq:sl-parseval}
unchanged.  This is the one-channel realization of the direct integral in
Section~\ref{subsec:direct-integrals}.

The self-adjoint expansion contains bound states and the physical continuum,
but no off-axis resonance projection.  Resonance poles appear only after a
boundary value of the Green kernel or $m$ function has been analytically
continued.  This clean separation prevents a Gamow vector from being inserted
as an extra orthogonal term in the physical spectral theorem.

\subsection{Volterra construction of the Jost solution}
\label{subsec:jost-volterra}
% BEGIN CURATED SUBJECT INDEX
\index{Jost function}
\index{Wronskian}
\index{boundary condition!incoming and outgoing}
\index{Riccati equation}
% END CURATED SUBJECT INDEX
\index{Volterra equation}\index{Jost solution}

Consider the $s$-wave equation in dimensionless form
\begin{equation}
  u''(r)+[k^2-U(r)]u(r)=0.
  \label{eq:s-wave-jost-equation}
\end{equation}
The Jost solution is fixed at infinity by
\begin{equation}
  f(k,r)\sim\rme^{\rmi kr}.
  \label{eq:s-wave-jost-asymptotic}
\end{equation}
Variation of constants gives the Volterra equation
\begin{equation}
  f(k,r)
  =\rme^{\rmi kr}
  +\frac{1}{k}\int_r^\infty
  \sin[k(r'-r)]U(r')f(k,r')\dd r'.
  \label{eq:jost-volterra}
\end{equation}
Applying $\rmd^2/\rmd r^2+k^2$ verifies the differential equation; the
upper-triangular integration region enforces the outgoing normalization.
Successive substitution produces a Neumann series.  Weighted integrability
of $U$ controls convergence and analyticity in $k$ in the appropriate half
plane.  For $l>0$, Riccati--Hankel functions replace the exponential and the
kernel contains the corresponding free radial Green function.

The regular solution follows from the forward Volterra equation
\begin{equation}
  \varphi(k,r)
  =\frac{\sin kr}{k}
  +\frac{1}{k}\int_0^r
  \sin[k(r-r')]U(r')\varphi(k,r')\dd r',
  \label{eq:regular-volterra}
\end{equation}
with Eq.~\eqref{eq:regular-solution-initial}.  The Jost function is the
constant Wronskian
\begin{equation}
  F(k)=W[f(k,\cdot),\varphi(k,\cdot)].
  \label{eq:sl-jost-wronskian-definition}
\end{equation}
At $r=0$, this convention gives $F(k)=f(k,0)$ for the $s$ wave.  At infinity,
the same Wronskian is the incoming coefficient of the regular solution, up to
the fixed normalization displayed in
Eq.~\eqref{eq:regular-jost-expansion}.  Constancy of the Wronskian is the
reason the origin domain and the radiation condition can be tested at
different ends of the interval.

\subsection{Worked example: the attractive spherical square well}
\label{subsec:square-well-jost}
% BEGIN CURATED SUBJECT INDEX
\index{Jost function}
\index{boundary condition!incoming and outgoing}
\index{Riemann sheet}
\index{branch cut}
\index{resonance!residue}
% END CURATED SUBJECT INDEX
\index{square-well potential}

Take
\begin{equation}
  U(r)=
  \begin{cases}
    -U_0,&0<r<a,\\
    0,&r>a,
  \end{cases}
  \qquad U_0>0,
  \label{eq:square-well-potential}
\end{equation}
and define $q=(k^2+U_0)^{1/2}$ on a stated branch.  The regular solution is
\begin{equation}
  \varphi(k,r)=\frac{\sin qr}{q},
  \qquad 0<r<a.
  \label{eq:square-well-regular-inside}
\end{equation}
For $r>a$ write
\begin{equation}
  \varphi(k,r)=A(k)\rme^{\rmi kr}+B(k)\rme^{-\rmi kr}.
  \label{eq:square-well-outside}
\end{equation}
Continuity at $a$ gives
\begin{align}
  2B(k)\rme^{-\rmi ka}
  &=\frac{\sin qa}{q}+\frac{\rmi}{k}\cos qa,
  \label{eq:square-well-incoming-coefficient}\\
  2A(k)\rme^{\rmi ka}
  &=\frac{\sin qa}{q}-\frac{\rmi}{k}\cos qa.
  \label{eq:square-well-outgoing-coefficient}
\end{align}
A convenient Jost normalization is therefore
\begin{equation}
  F(k)=\rme^{\rmi ka}
  \left[\cos qa-\rmi\frac{k}{q}\sin qa\right].
  \label{eq:square-well-jost-function}
\end{equation}
The purely outgoing condition $B(k)=0$ becomes
\begin{equation}
  q\cot(qa)=\rmi k.
  \label{eq:square-well-resonance-condition}
\end{equation}
For $k=\rmi\kappa$ on the physical sheet it gives a decaying bound state;
continued into the lower half-plane it gives virtual and resonant poles.
No new differential equation was introduced when the pole changed its name:
only the sheet and asymptotic behavior changed.

To compute a residue, differentiate Eq.~\eqref{eq:square-well-jost-function}
on the same branches of $k$ and $q$.  The result enters as
$F'(k_{\mathrm R})$ in Eq.~\eqref{eq:radial-green-residue}.  Numerically,
finite differencing $F$ across a branch cut will give the wrong normalization
even if a root finder happened to locate the pole.

\subsection{Wronskian derivatives and Zel'dovich normalization}
% BEGIN CURATED SUBJECT INDEX
\index{GNS construction}
\index{Jost function}
\index{Wronskian}
\index{exact WKB}
\index{monodromy}
\index{Zel'dovich regularization}
\index{Berggren completeness relation}
% END CURATED SUBJECT INDEX
\index{Zel'dovich regularization}\index{Jost-function derivative}
\label{subsec:wronskian-derivative-normalization}

The relation between a Wronskian derivative and a bilinear norm follows from
Lagrange's identity with two nearby spectral parameters.  For simplicity let
\begin{equation}
  [-\partial_r^2+U(r)-k^2]u(k,r)=0.
  \label{eq:parameter-dependent-radial}
\end{equation}
Differentiate with respect to $k$ and subtract the equation multiplied in the
opposite order.  Integration to $R$ gives
\begin{equation}
  2k\int_0^R u(k,r)^2\dd r
  =W[u,\partial_k u](R)-W[u,\partial_k u](0).
  \label{eq:wronskian-derivative-integral}
\end{equation}
There is no complex conjugation: the resonance normalization is bilinear.
The origin term is fixed by the regular normalization.  For a bound state the
term at infinity vanishes after the decaying solution is normalized.  For a
resonance it diverges on the real axis, but analytic continuation or a
Zel'dovich Gaussian factor assigns the continued boundary value.  The result
is proportional to $F'(k_{\mathrm R})$.

Equation~\eqref{eq:wronskian-derivative-integral} is the calculational core of
Berggren's normalization proof in Section~\ref{subsec:berggren-proof}: the
regularized integral is not postulated independently of the pole.  It is the
continued derivative of the same Wronskian that defines the pole.  Complex
scaling makes the boundary term vanish along a rotated ray; exact WKB
calculates it by connection and monodromy data.  These are different
realizations of one analytic identity.

\subsection{Exact WKB computes the Jost Wronskian globally}
\label{subsec:sl-jost-exact-wkb}
% BEGIN CURATED SUBJECT INDEX
\index{Hilbert space}
\index{Jost function}
\index{Wronskian}
\index{connection coefficient}
\index{boundary condition!incoming and outgoing}
\index{resonance!residue}
\index{exact WKB}
\index{Langer correction}
\index{Frobenius solution}
\index{complex scaling}
% END CURATED SUBJECT INDEX

Let $\vsub{\psi}{reg}$ denote the exact-WKB solution selected by the origin
domain, including the Frobenius and Langer data when the endpoint is singular.
Let $\vsub{\psi}{out}$ be the Borel-resummed branch selected by the outgoing
sector at infinity.  Then
\begin{equation}
  \Cincoming(E,\hbar)
  \propto
  W[\vsub{\psi}{out},\vsub{\psi}{reg}]
  \propto F_l(k(E)).
  \label{eq:wkb-jost-wronskian}
\end{equation}
Local WKB series do not determine this Wronskian.  One must transport the two
solutions through the Stokes graph, multiply the connection matrices in the
correct order, and retain the exponentially small entries.  A zero of the
result removes the incoming component.  Its derivative gives the residue by
Eqs.~\eqref{eq:radial-green-residue} and
\eqref{eq:wronskian-derivative-integral}.

This is the novel reorganization used in the rest of the book.  Sturm--
Liouville theory determines the physical operator and continuum measure;
Jost theory identifies the analytic Wronskian; exact WKB calculates that
Wronskian nonperturbatively.  Complex scaling and rigged Hilbert spaces will
change the realization of the resulting pole but not the connection zero or
its derivative.

\subsection{Where Coulomb and many-body asymptotics change the construction}
% BEGIN CURATED SUBJECT INDEX
\index{operator!adjoint and self-adjoint}
\index{spectral theorem}
\index{Jost function}
\index{Wronskian}
\index{scattering!Coulomb}
\index{scattering!three-body Coulomb}
\index{infrared problem}
\index{dressing state}
% END CURATED SUBJECT INDEX
\index{Coulomb scattering}\index{three-body Coulomb scattering}
\label{subsec:sl-jost-long-range-limit}

For a two-body Coulomb tail, the half-line differential operator and its
self-adjoint spectral theorem remain valid, but the comparison solutions are
not $\rme^{\pm\rmi kr}$.  One uses Coulomb--Hankel functions with the
logarithmic phase
\begin{equation}
  kr-\eta\log(2kr)-\frac{l\pi}{2}+\sigma_l.
  \label{eq:coulomb-jost-phase}
\end{equation}
The Coulomb-modified Jost function is a Wronskian with those reference
solutions.  Applying Eq.~\eqref{eq:jost-volterra} with the free kernel to the
full Coulomb potential fails because the integral is not short ranged.

For three charged particles, different cluster regions carry different
logarithmic phases and no single universally usable scalar asymptotic Jost
basis is known.  A finite-dimensional coupled equation can still be solved in
a controlled interaction region, but the global comparison dynamics must be
declared separately.  This is the same structural difficulty that appears as
infrared dressing in QED: long-range correlations change the asymptotic state
space before one asks for a pole.  Section~\ref{sec:long-range-ir} develops
that obstruction rather than hiding it in a formal plane wave.

\subsection{A reproducible half-line pole calculation}
\label{subsec:sl-jost-algorithm}
% BEGIN CURATED SUBJECT INDEX
\index{operator!adjoint and self-adjoint}
\index{resolvent}
\index{Sturm--Liouville theory}
\index{Sturm--Liouville theory!limit-point/limit-circle}
\index{Jost function}
\index{Wronskian}
\index{boundary condition!incoming and outgoing}
\index{resonance!residue}
\index{scattering!Coulomb}
% END CURATED SUBJECT INDEX

For later laboratories, the complete calculation can be organized into the
following steps:
\begin{enumerate}[label=\arabic*.]
  \item specify $p,q,w$ and the Hilbert-space measure;
  \item classify every singular endpoint as limit point or limit circle;
  \item impose only the boundary conditions required for a self-adjoint
  physical Hamiltonian;
  \item construct the regular solution by an initial-value or Volterra
  equation;
  \item choose the correct short-range or Coulomb comparison solution at
  infinity and construct the Jost solution;
  \item form their Wronskian and verify that it is independent of the matching
  point;
  \item declare the momentum sheet before continuing and finding a zero;
  \item calculate the Wronskian derivative and a tested resolvent residue;
  \item vary the matching point, integration contour, and numerical basis
  while tracking the pole as one analytic object.
\end{enumerate}

\begin{checkpointbox}
The reader should now be able to derive the Green kernel from its derivative
jump, explain why an $s$-wave origin condition is domain data rather than a
numerical convenience, obtain the square-well condition
Eq.~\eqref{eq:square-well-resonance-condition}, and identify the exact-WKB
incoming coefficient with a continued Sturm--Liouville Wronskian.
\end{checkpointbox}
